\documentclass[11pt]{amsart}

\usepackage[T1]{fontenc}
\usepackage{lmodern}
\usepackage{microtype}
\usepackage{amsmath,amssymb,amsthm}
\usepackage[hidelinks]{hyperref}

\hypersetup{
  pdftitle={Counterexamples to the Hypercube Maximum Propagation-Time Conjecture},
  pdfsubject={Zero forcing and propagation time on hypercubes},
  pdfkeywords={zero forcing, propagation time, hypercube}
}

\newtheorem{theorem}{Theorem}
\theoremstyle{remark}
\newtheorem*{remark}{Remark}

\DeclareMathOperator{\pt}{pt}
\DeclareMathOperator{\PT}{PT}
\DeclareMathOperator{\Z}{Z}

\title{Counterexamples to the Hypercube Maximum Propagation-Time Conjecture}
\date{}
\subjclass[2020]{05C50, 05C57}
\keywords{zero forcing, propagation time, hypercube}

\begin{document}

\begin{abstract}
Brimkov, Cameron, and Grubbs conjectured that the maximum propagation
time of the hypercube $Q_d$, taken over its minimum zero forcing sets,
is $2^{d-2}$. We give explicit minimum zero forcing sets of $Q_6$ and
$Q_7$ with propagation times $18$ and $43$, respectively. Hence
$\PT(Q_6)\ge 18>16$ and $\PT(Q_7)\ge 43>32$, disproving the conjectured
formula. Cameron and Pulaj reported the conjectured values for
$2\le d\le5$, so $Q_6$ is the smallest-dimensional counterexample to
this formula. We do not determine either exact maximum or address the
full propagation-time-interval assertion in the original conjecture.
\end{abstract}

\maketitle

\section{The counterexamples}

We use the standard synchronous zero forcing convention
of~\cite{HHKMWY}. Identify the vertices of $Q_d$ with
$V_d=\{0,1,\ldots,2^d-1\}$ through their length-$d$ binary expansions.
Thus $x,y\in V_d$ are adjacent exactly when $x\mathbin{\oplus}y$ is a
power of two, where $\oplus$ denotes bitwise exclusive-or. For
$C\subseteq V_d$, let
\[
\Phi_d(C)=
\bigl\{v\in V_d\setminus C:
  \text{there is }u\in C\text{ with }N(u)\setminus C=\{v\}\bigr\},
\]
where $N(u)$ is the open neighborhood of $u$. Starting from $C_0=S$,
define
\[
D_t=\Phi_d(C_{t-1}),
\qquad
C_t=C_{t-1}\cup D_t
\quad(t\ge1).
\]
Choosing one witness $u$ for each $v\in D_t$ gives independent forces,
since a vertex with a unique neighbor outside $C_{t-1}$ cannot witness
two distinct targets. Thus these are exactly the synchronous forcing
layers. If $C_T=V_d$ and $C_{T-1}\ne V_d$, then
$\pt(Q_d,S)=T$. The parameter $\PT(G)$ is the maximum of $\pt(G,S)$
over all minimum zero forcing sets $S$ of $G$.

Brimkov, Cameron, and Grubbs conjectured that
\[
\PT(Q_d)=2^{d-2}\qquad(d\ge2)
\]
as part of their Conjecture~4.4~\cite{BCG}. Cameron and Pulaj restated
this conjecture and reported its validity for $2\le d\le5$
\cite[Conjecture~2.1 and Table~7]{CP}.

\begin{theorem}\label{thm:main}
Let
\begin{align*}
S_6=\{&2,5,6,7,10,13,16,17,18,20,21,22,25,28,29,30,31,\\
      &33,35,36,38,43,44,49,50,51,56,57,58,60,62,63\}
\end{align*}
and
\begin{align*}
S_7=\{&0,1,2,3,4,10,11,13,14,15,21,22,24,27,30,31,\\
      &34,36,37,38,39,41,45,46,48,49,50,51,55,56,57,58,63,\\
      &64,65,66,67,68,69,71,74,77,79,85,88,90,93,98,100,101,\\
      &102,103,108,109,112,113,115,119,120,122,123,124,126,127\}.
\end{align*}
Then $S_6$ and $S_7$ are minimum zero forcing sets and
\[
\pt(Q_6,S_6)=18,
\qquad
\pt(Q_7,S_7)=43.
\]
Consequently,
\[
\PT(Q_6)\ge18>2^{6-2},
\qquad
\PT(Q_7)\ge43>2^{7-2}.
\]
In particular, the maximum-propagation-time assertion in
\cite[Conjecture~4.4]{BCG} is false.
\end{theorem}

\begin{proof}
The known identity $\Z(Q_d)=2^{d-1}$
\cite[Theorem~3.1]{AIM} gives $\Z(Q_6)=32$ and $\Z(Q_7)=64$.
The displayed sets have these respective cardinalities. It therefore
suffices to verify that they force all vertices and to compute their
propagation times. All calculations below use
\[
N(u)=\{u\mathbin{\oplus}2^i:0\le i<d\}
\]
and exact set operations.

For $S_6$, direct evaluation of the recurrence gives
\begin{align*}
(D_1,\ldots,D_6)
 &= (\{61\},\{52\},\{4\},\{14\},\{26\},\{19\}),\\
(D_7,\ldots,D_{11})
 &= (\{1\},\{37\},\{32\},\{41\},\{59\}),\\
(D_{12},\ldots,D_{15})
 &= (\{27,42,55\},\{11,24,47\},\{8,9,12\},\{0\}),\\
(D_{16},D_{17},D_{18})
 &= (\{3,40,48\},\{15,23,34,45,53\},\{39,46,54\}).
\end{align*}
The listed layers are disjoint and their union is $V_6\setminus S_6$.
Hence $C_{18}=V_6$ and $D_{18}\ne\varnothing$, so
$\pt(Q_6,S_6)=18$.

For $S_7$, the exact layers are
\begin{align*}
(D_1,\ldots,D_{10})
 &= (\{5\},\{7\},\{47\},\{35\},\{19\},\{59\},\{43\},
     \{111\},\{99\},\{33\}),\\
(D_{11},\ldots,D_{20})
 &= (\{53\},\{17\},\{9\},\{75\},\{83\},\{114\},\{106\},
     \{96\},\{116\},\{80\}),\\
(D_{21},\ldots,D_{30})
 &= (\{72\},\{78\},\{95\},\{125\},\{105\},\{40\},\{60\},
     \{92\},\{89\},\{29\}),\\
(D_{31},\ldots,D_{36})
 &= (\{12,23\},\{6,18,20,54,87\},\{70,86\},\{82\},
     \{118\},\{110,117\}),\\
(D_{37},\ldots,D_{40})
 &= (\{97,107\},\{32,42\},\{104\},\{121\}),\\
(D_{41},D_{42},D_{43})
 &= (\{61,73,81,91\},\{25,44,52,62,76,84,94\},
     \{8,16,26,28\}).
\end{align*}
These layers are disjoint and their union is $V_7\setminus S_7$.
Thus $C_{43}=V_7$ and $D_{43}\ne\varnothing$, proving
$\pt(Q_7,S_7)=43$.

Both sets are therefore zero forcing sets of minimum cardinality. The
bounds on $\PT(Q_6)$ and $\PT(Q_7)$ follow from its definition.
\end{proof}

\begin{remark}
The $Q_6$ certificate alone disproves the formula
$\PT(Q_d)=2^{d-2}$. The theorem gives only lower bounds on
$\PT(Q_6)$ and $\PT(Q_7)$; it does not determine either exact maximum.
It also does not decide whether $Q_6$ or $Q_7$ has a full realized
propagation-time interval, the separate interval assertion in
\cite[Conjecture~4.4]{BCG}.
\end{remark}

\begin{thebibliography}{9}

\bibitem{HHKMWY}
L.~Hogben, M.~Huynh, N.~Kingsley, S.~Meyer, S.~Walker, and M.~Young,
\emph{Propagation time for zero forcing on a graph},
Discrete Appl. Math. \textbf{160} (2012), no.~13--14, 1994--2005,
\href{https://doi.org/10.1016/j.dam.2012.04.003}
{doi:10.1016/j.dam.2012.04.003}.

\bibitem{BCG}
B.~Brimkov, T.~R. Cameron, and O.~Grubbs,
\emph{On the forts and related parameters of the hypercube graph},
\href{https://arxiv.org/abs/2507.10826}{arXiv:2507.10826v1} (2025).

\bibitem{CP}
T.~R. Cameron and J.~Pulaj,
\emph{IP models for minimum zero forcing sets, forts, and related graph
parameters},
\href{https://arxiv.org/abs/2508.07293}{arXiv:2508.07293v1} (2025).

\bibitem{AIM}
AIM Minimum Rank--Special Graphs Work Group,
\emph{Zero forcing sets and the minimum rank of graphs},
Linear Algebra Appl. \textbf{428} (2008), no.~7, 1628--1648,
\href{https://doi.org/10.1016/j.laa.2007.10.009}
{doi:10.1016/j.laa.2007.10.009}.

\end{thebibliography}

\end{document}